\chapter{Algebraic surgery, the exponential machine, and the squeeze}

\begin{orient}
\textbf{What this chapter is for.} Three whole families of limit yield without a single derivative, a single series, and a single appeal to L'H\^opital. The first family is \emph{surgery on the expression}: conjugates, factorisation of a removable singularity, common denominators, and the deliberate insertion of a bridging term, all of which turn an indeterminate shape into a determinate one by pure algebra. The second is the \emph{exponential machine}: every power-form indeterminacy ($1^{\infty}$, $0^{0}$, $\infty^{0}$) becomes a product the moment you write $f^{g}=\exp(g\ln f)$, and the whole $e$-family falls out as one corollary. The third is the \emph{squeeze}, which is the only tool in the book that works when the limit exists but nothing inside the expression converges at all. By the end you should be able to look at an indeterminate form and say, in a few seconds, which of the three families it belongs to, and you should be able to execute all three without hesitation and without reaching for calculus you do not need.
\end{orient}

\section{Choosing the family}

An indeterminate form is not one problem but a symptom shared by several. Before choosing a technique, ask what is actually producing the indeterminacy. If it is a \emph{cancellation} that the algebra is hiding, whether a common factor, a difference of radicals, or two fractions that must be combined, then the cure is surgery and nothing else is needed. If it is an \emph{exponent} fighting a base, the cure is the logarithm, which demotes the fight to a product. If the expression contains a piece that simply does not converge, such as $\sin(1/x)$ near $0$, or a fractional part, or a maximum over a growing index set, then no amount of rewriting will produce a limit for that piece, and the only honest move is to bound the whole expression from both sides.

The three questions are worth asking in that order, and the order is not arbitrary. Surgery is cheap and often removes the problem entirely, leaving a value you can read off by substitution. The logarithm is a mechanical transform that never fails but always costs you a second limit, so it is worth paying for only when the first question has been answered ``no''. The squeeze demands that you invent bounds, which is the only genuinely creative step of the three, and inventing bounds is work you should not do until you know the other two will not serve.

It also helps to name the form precisely rather than calling everything ``indeterminate''. There are seven: $\tfrac00$, $\tfrac{\infty}{\infty}$, $0\cdot\infty$, $\infty-\infty$, $1^{\infty}$, $0^{0}$, and $\infty^{0}$. The last three are all powers and all go to the same place, the logarithm. The middle two are one algebraic step from a quotient, since $0\cdot\infty$ becomes $\tfrac00$ on inverting either factor and $\infty-\infty$ becomes a quotient on taking a common denominator or a conjugate. So the seven forms are really three: a quotient, a power, and a nonconvergent piece. That is the whole taxonomy this chapter needs.

There is a fourth answer that this chapter deliberately postpones. If the expression is a sum of $n$ terms with a $1/n$ in front, or a difference quotient in disguise, the right reading is neither algebraic nor exponential nor a squeeze: it is an integral or a derivative, and that is the subject of the next chapter. The decision tree below sends you there when none of the three families fits.

\begin{center}
\begin{tikzpicture}[node distance=6mm and 8mm]
\node[dtest] (a) {Is the expression a power $f^{g}$\\ with the exponent varying?};
\node[dnode, below left=9mm and 6mm of a] (b) {Logarithm machine:\\ $\exp\bigl(\lim g\ln f\bigr)$};
\node[dtest, below right=9mm and 2mm of a] (c) {Do radicals, a common factor,\\ or split fractions cause it?};
\node[dleaf, below left=9mm and 1mm of c] (d) {Algebraic surgery:\\ conjugate, factor, bridge};
\node[dtest, below right=9mm and 1mm of c] (e) {Is some factor bounded\\ but non-convergent?};
\node[dleaf, below left=9mm and 0mm of e] (f) {Squeeze; bounded\\ times infinitesimal};
\node[dnode, below right=9mm and 0mm of e] (g) {Expansion, or read it\\ as an integral};
\draw[dedge] (a) -- node[dlab,left]{yes} (b);
\draw[dedge] (a) -- node[dlab,right]{no} (c);
\draw[dedge] (c) -- node[dlab,left]{yes} (d);
\draw[dedge] (c) -- node[dlab,right]{no} (e);
\draw[dedge] (e) -- node[dlab,left]{yes} (f);
\draw[dedge] (e) -- node[dlab,right]{no} (g);
\end{tikzpicture}
\end{center}

\section{Surgery on the expression}

Every technique in this section rests on one observation: an indeterminate form is a statement about the \emph{written expression}, not about the function. The function $\dfrac{x^{2}-4}{x-2}$ is the function $x+2$ with one point deleted, and the $\tfrac00$ at $x=2$ is an artefact of how somebody chose to write it down. Surgery is the art of rewriting until the artefact is gone. The five moves below are the complete toolkit for the rational and radical cases, and they are ordered by how visible the hidden cancellation is: a common polynomial factor is the most visible, a bridging term the least.

It is worth being explicit about what ``determinate'' means here, because that is the target of every rewrite. A quotient is determinate at $a$ when the denominator does not vanish there and both parts are continuous, so that substitution is legal. Every move in this section aims at that state, and you should stop rewriting the instant you reach it. The most common failure of technique in this whole chapter is continuing to manipulate after the problem has already been solved.

\begin{keynote}[The standard limits these moves reduce to]
Surgery is only useful because a small list of limits is known outright. Memorise these six and you will never need a series expansion for a first-order question:
\[\lim_{x\to0}\frac{\sin x}{x}=1,\quad \lim_{x\to0}\frac{1-\cos x}{x^{2}}=\frac12,\quad \lim_{x\to0}\frac{\eu^{x}-1}{x}=1,\]
\[\lim_{x\to0}\frac{\ln(1+x)}{x}=1,\quad \lim_{x\to0}\frac{a^{x}-1}{x}=\ln a,\quad \lim_{x\to0}\frac{(1+x)^{\alpha}-1}{x}=\alpha .\]
The last of these is what a conjugate computation reproduces by hand for $\alpha=\tfrac12$ and $\alpha=\tfrac13$; the conjugate is the elementary proof, not a shortcut past it.
\end{keynote}

\begin{keynote}[When surgery is the wrong choice]
Surgery fails, and should be abandoned quickly, when the cancelling factor is transcendental rather than algebraic. In $\dfrac{\eu^{x}-1-x}{x^{2}}$ there is no factor to divide out and no conjugate to apply: the numerator is not a difference of powers of anything. That limit is $\tfrac12$, and it belongs to the expansion techniques, not here. The diagnostic is simple. If every vanishing piece is a polynomial or a root of one, surgery will finish the job; if any vanishing piece is $\eu^{x}$, $\ln$, or a trigonometric function set against its own linearisation, it will not.
\end{keynote}

Start with radicals, because a difference of radicals is the most common way to disguise a cancellation. The mechanism is always the same: a difference of $n$th roots is hard, but a difference of $n$th \emph{powers} is easy, so multiply by whatever converts one into the other.

\tech{Conjugate multiplication for a difference of square roots}{easy}
\cue Any $\sqrt{A}-\sqrt{B}$ in a numerator or a denominator, where $A-B$ is visibly simpler than either radical. The classic signatures are $\sqrt{x+h}-\sqrt{x}$, $\sqrt{x^{2}+ax}-x$, and a denominator that vanishes because two radicals become equal in the limit.
\method Multiply the whole fraction by $\dfrac{\sqrt A+\sqrt B}{\sqrt A+\sqrt B}$. The identity is
\[\sqrt{A}-\sqrt{B}=\frac{A-B}{\sqrt{A}+\sqrt{B}},\]
which trades a difference that is $0-0$ for a quotient whose numerator carries the cancelling factor explicitly and whose denominator tends to a harmless nonzero constant. If the radicals sit in the denominator, the same multiplication rationalises it and exposes the factor that must be cancelled from the numerator.

\begin{worked}
\[\lim_{x\to0}\frac{\sin^{2}x}{\sqrt{2}-\sqrt{1+\cos x}}.\]
Both parts vanish. Multiply top and bottom by $\sqrt2+\sqrt{1+\cos x}$; the denominator becomes $2-(1+\cos x)=1-\cos x$, so the quotient is
\[\frac{\sin^{2}x\,\bigl(\sqrt2+\sqrt{1+\cos x}\bigr)}{1-\cos x}=(1+\cos x)\bigl(\sqrt2+\sqrt{1+\cos x}\bigr),\]
using $\sin^{2}x=(1-\cos x)(1+\cos x)$ to cancel the vanishing factor outright. Now substitute: the value is $2\cdot 2\sqrt2=4\sqrt2\approx5.65685$. No series, no derivatives, and note that the notorious $\tfrac{1-\cos x}{x^{2}}$ never had to be invoked.
\end{worked}

The same identity handles $\infty-\infty$ at infinity, where it is the standard first move and where the arithmetic is if anything easier. Three specimens worth carrying:
\[\sqrt{x^{2}+x+1}-\sqrt{x^{2}-x+1}=\frac{2x}{\sqrt{x^{2}+x+1}+\sqrt{x^{2}-x+1}}\longrightarrow1,\]
because the denominator is $2x\bigl(1+o(1)\bigr)$; likewise $x\bigl(\sqrt{x^{2}+1}-x\bigr)=\dfrac{x}{\sqrt{x^{2}+1}+x}\to\tfrac12$, and $\sqrt{x^{2}+ax}-x\to\tfrac a2$. In each case the conjugate converts a cancellation of leading terms, which floating-point arithmetic also finds catastrophic, into an honest quotient.

\begin{trap}
Multiplying only the denominator by its conjugate. You must multiply the \emph{whole fraction} by $\tfrac{\text{conjugate}}{\text{conjugate}}$, which is multiplication by $1$; touching one storey only changes the function. The second common error occurs at $x\to-\infty$, where $\sqrt{x^{2}}=\abs{x}=-x$, so factoring $x$ out of a radical carries a sign that is easy to lose and that flips the answer.
\end{trap}

\begin{seealsob}The cube-root conjugate; add and subtract a bridging term; change of variable to move the limit point.\end{seealsob}

Square roots are the easy case only because $a^{2}-b^{2}$ factors in one step. For cube roots the same idea works, but the multiplier is different, and guessing it is where most attempts die.

\tech{The cube-root conjugate}{medium}
\cue A difference $\sqrt[3]{A}-\sqrt[3]{B}$, or a cube root minus a constant, inside a $\tfrac00$ quotient. Also any mixed $\sqrt[3]{A}-\sqrt{B}$, which is handled by first splitting it with a bridging term and then applying this move to the cube-root half.
\method Set $a=\sqrt[3]{A}$, $b=\sqrt[3]{B}$ and use
\[a-b=\frac{a^{3}-b^{3}}{a^{2}+ab+b^{2}}=\frac{A-B}{a^{2}+ab+b^{2}}.\]
The denominator tends to $3L^{2}$ when both roots tend to a common nonzero $L$, so it is inert and all the content sits in $A-B$. The general $n$th-root version is
\[a-b=\frac{a^{n}-b^{n}}{a^{n-1}+a^{n-2}b+\cdots+ab^{n-2}+b^{n-1}},\]
with $n$ terms in the denominator, each tending to $L^{n-1}$.

\begin{worked}
\[\lim_{x\to0}\frac{\sqrt[3]{8+3x-x^{2}}-2}{x+x^{2}}.\]
With $a=\sqrt[3]{8+3x-x^{2}}$ and $b=2$ we have $a^{3}-b^{3}=3x-x^{2}=x(3-x)$, while $a^{2}+2a+4\to 4+4+4=12$. Hence the quotient equals
\[\frac{x(3-x)}{(x+x^{2})\bigl(a^{2}+2a+4\bigr)}=\frac{3-x}{(1+x)\bigl(a^{2}+2a+4\bigr)}\longrightarrow\frac{3}{12}=\frac14 .\]
A mixed instance, resolved by the same identity applied twice:
\[\lim_{x\to0}\frac{\sqrt[3]{1+x}-\sqrt{1+x}}{x}=\frac13-\frac12=-\frac16\approx-0.166667,\]
where each root is compared with $1$ separately, the cube root by the identity above and the square root by the previous technique.
\end{worked}

\begin{trap}
Guessing the wrong multiplier. The cube conjugate of $a-b$ is $a^{2}+ab+b^{2}$, with all plus signs; $a^{2}-ab+b^{2}$ is the multiplier for $a+b$. Writing the wrong one produces an expression that is simply not equal to the original, and the error is silent because the answer still looks plausible. Check the multiplier by expanding it against $a-b$ before you use it.
\end{trap}

\begin{seealsob}Conjugate multiplication for a difference of square roots; add and subtract a bridging term; change of variable to move the limit point.\end{seealsob}

When no radicals are present, the hidden cancellation is a common polynomial factor, and finding it is a division problem rather than an identity problem. This is the most elementary technique in the chapter and also the one most often skipped in favour of machinery.

\tech{Factor out the removable singularity}{easy}
\cue A rational expression, or one that factors, evaluating to $\tfrac00$ at $x=a$. Equally: a difference of two fractions, each of which blows up at $a$, whose combination over a common denominator is finite.
\method If numerator and denominator both vanish at $x=a$, then $(x-a)$ divides both, by the factor theorem. Divide it out, by inspection or by synthetic division, and repeat until at least one of the two no longer vanishes at $a$. Then substitute. When the expression is a difference of fractions, combine over a common denominator \emph{first}; the cancellation only becomes visible once there is a single fraction to look at.

\begin{worked}
\[\lim_{x\to2}\frac{x^{4}-6x^{3}+13x^{2}-12x+4}{x^{3}-5x^{2}+8x-4}.\]
Both vanish at $x=2$, and both still vanish after one division, so $x=2$ is at least a double root of each. In fact the numerator is $(x-1)^{2}(x-2)^{2}$ and the denominator is $(x-1)(x-2)^{2}$, so on a punctured neighbourhood the quotient is exactly $x-1$, and the limit is $1$.

A second shape, the split fraction, where the common denominator is the whole trick:
\[\lim_{x\to1}\left(\frac{1}{1-x}-\frac{3}{1-x^{3}}\right)=\lim_{x\to1}\frac{(1+x+x^{2})-3}{1-x^{3}}=\lim_{x\to1}\frac{(x-1)(x+2)}{-(x-1)(1+x+x^{2})}=-1 .\]
Neither fraction has a limit at $x=1$; their difference does, and only the combined form shows it.

A third, worth knowing as a formula: $\displaystyle\lim_{x\to1}\frac{x^{m}-1}{x^{n}-1}=\frac{m}{n}$, since both factor as $(x-1)$ times a sum of $m$ or $n$ terms each tending to $1$.
\end{worked}

\begin{trap}
Reaching for L'H\^opital. It gives the same answer here, but it destroys the structural information: the point of the calculation is that the singularity was \emph{removable} and the function extends continuously across $x=a$, a fact that continuity arguments and radius-of-convergence arguments later depend on. The second error is stopping after one division when the root is repeated, as above, and concluding that the limit is still indeterminate.
\end{trap}

\begin{seealsob}The cube-root conjugate; change of variable to move the limit point; the limit is a derivative.\end{seealsob}

The three moves so far all require a shared structure to exploit: a conjugate needs matching radicals, a factorisation needs a common root. What if there is none, because the two competing pieces have nothing algebraically in common? Then you manufacture the structure, by splitting the difference at its common limit value. This is telescoping applied to a difference of two terms rather than to a long sum, and it is the reason the technique looks like an accounting trick rather than an identity.

\tech{Add and subtract a bridging term}{medium}
\cue A difference of two dissimilar expressions that tend to the same value: $\sqrt[3]{f(x)}-\sqrt{g(x)}$ with $f,g\to1$, or $\eu^{u(x)}-\cos v(x)$, or any $A-B$ where $A$ and $B$ are each tractable against a constant but not against each other.
\method Insert the common limit $L$ and telescope the difference:
\[A-B=(A-L)-(B-L).\]
Each bracket now has a single radical or a single transcendental function set against a constant, so each yields to its own conjugate or its own standard limit. Take the two limits separately and subtract. The same device extends to three or more pieces, bridging at $L$ each time.

\begin{worked}
\[\lim_{x\to0}\frac{\sqrt[3]{1+\sin^{2}x}-\sqrt{\cos x}}{x^{2}}.\]
Both radicals tend to $1$, so bridge at $L=1$:
\[\frac{\sqrt[3]{1+\sin^{2}x}-1}{x^{2}}+\frac{1-\sqrt{\cos x}}{x^{2}}.\]
For the first, the cube conjugate gives $\dfrac{\sin^{2}x}{x^{2}\bigl(a^{2}+a+1\bigr)}\to\dfrac{1}{3}$ with $a=\sqrt[3]{1+\sin^{2}x}\to1$. For the second, the square conjugate gives $\dfrac{1-\cos x}{x^{2}\bigl(1+\sqrt{\cos x}\bigr)}\to\dfrac{1/2}{2}=\dfrac14$. The limit is $\tfrac13+\tfrac14=\tfrac{7}{12}\approx0.583333$.

The same split on a cleaner pair confirms the bookkeeping: $\dfrac{\sqrt[3]{1+3x}-\sqrt{1+2x}}{x^{2}}\to-\tfrac12$, since $\sqrt[3]{1+3x}=1+x-x^{2}+\cdots$ and $\sqrt{1+2x}=1+x-\tfrac{x^{2}}{2}+\cdots$, so the linear parts cancel and the quadratic parts do not.
\end{worked}

\begin{trap}
Splitting a limit into two pieces that individually diverge. The identity $\lim(P+Q)=\lim P+\lim Q$ requires both right-hand limits to exist finitely. Bridging at the common value $L$ is exactly the device that guarantees this, which is why $L$ must be the shared limit of $A$ and $B$ and not merely some convenient constant. If $A$ and $B$ tend to different values, the original difference is not indeterminate and there is nothing to bridge.
\end{trap}

\begin{seealsob}Conjugate multiplication for a difference of square roots; the cube-root conjugate; log-transform a product or a nested power.\end{seealsob}

All four moves are easier to see when the limit point is $0$. That is not a coincidence: every standard limit in the keynote above is written at $0$, and every series expansion in the next chapter is written at $0$. The last move of this section is the bookkeeping that puts you there, and it is also the move that turns competing radical orders into competing polynomial degrees.

\tech{Change of variable to move the limit point}{easy}
\cue The limit point is awkward ($x\to1$, $x\to\pi/2$, $x\to\infty$), or the expression is a function of a compound argument such as $\sqrt[12]{x}$ or $\ln x$, or several roots of different orders appear together.
\method Substitute so that the limit point becomes $0$ or $\infty$ and the roots become polynomial: $h=x-a$, $t=1/x$, $u=\pi/2-x$, or $x=t^{m}$ with $m=\lcm$ of the root orders. State the new direction of approach explicitly, because a substitution can convert a two-sided limit into a one-sided one and can reverse orientation.

\begin{worked}
\[\lim_{x\to1}\frac{\sqrt[3]{x}-1}{\sqrt[4]{x}-1}.\]
The orders are $3$ and $4$, so put $x=t^{12}$; as $x\to1$ through positive values, $t\to1$. The quotient becomes a ratio of polynomials with a visible common factor:
\[\frac{t^{4}-1}{t^{3}-1}=\frac{(t-1)(t^{3}+t^{2}+t+1)}{(t-1)(t^{2}+t+1)}\longrightarrow\frac{4}{3}.\]
A second instance, moving a trigonometric limit point to the origin: $\displaystyle\lim_{x\to\pi/2}\frac{\cos x}{x-\pi/2}$ becomes, under $u=x-\pi/2$,
\[\lim_{u\to0}\frac{\cos\bigl(\tfrac\pi2+u\bigr)}{u}=\lim_{u\to0}\frac{-\sin u}{u}=-1 .\]
A third, moving infinity to the origin: $\displaystyle\lim_{x\to\infty}x\ln\Bigl(1+\frac1x\Bigr)$ becomes $\displaystyle\lim_{t\to0^{+}}\frac{\ln(1+t)}{t}=1$ under $t=1/x$, where the standard limit applies directly.
\end{worked}

\begin{trap}
Substituting without transforming the direction of approach. Under $t=1/x$ the one-sided limit $x\to0^{-}$ becomes $t\to-\infty$, not $t\to+\infty$, and every sign in the answer depends on it. Under $x=t^{2}$ the variable $t$ ranges over one side only, so a genuinely two-sided limit in $x$ must either be checked for $t$ of both signs or the substitution restricted to $x>0$ with the other side treated separately.
\end{trap}

\begin{seealsob}Factor out the removable singularity; exponent simplification before base analysis; Riemann sum recognition.\end{seealsob}

The five moves are meant to be combined, and a problem of any weight will need two or three of them in sequence. A short illustration of the protocol in full:
\[\lim_{x\to0}\frac{\sqrt[3]{1+x}-\sqrt[3]{1-x}}{\sqrt{1+x}-\sqrt{1-x}}.\]
The form is $\tfrac00$. There is no common factor to divide out, so factorisation is not the move. Both storeys are differences of radicals, so apply the appropriate conjugate to each: the numerator becomes $\dfrac{2x}{a^{2}+ab+b^{2}}$ with $a=\sqrt[3]{1+x}$ and $b=\sqrt[3]{1-x}$, both tending to $1$, and the denominator becomes $\dfrac{2x}{\sqrt{1+x}+\sqrt{1-x}}$. The factors of $2x$ cancel, which is the removable singularity making itself visible, and what remains is
\[\frac{\sqrt{1+x}+\sqrt{1-x}}{a^{2}+ab+b^{2}}\longrightarrow\frac{2}{3}\approx0.666667 .\]
Three moves, no calculus, and the answer is exact. Note the order: conjugate first to expose the factor, cancel second, substitute third, and stop.

\section{The exponential machine}

A power $f(x)^{g(x)}$ with both $f$ and $g$ varying is not directly attackable: there is no arithmetic rule for the limit of a power when the base tends to $1$ and the exponent to $\infty$, because $1^{n}=1$ pulls one way and $c^{\infty}$ pulls the other, and both pulls are real. The resolution is not a special theorem but a change of representation. For $f>0$,
\[f^{g}=\exp\bigl(g\ln f\bigr),\]
and $\exp$ is continuous on all of $\R$. Continuity is the entire content of the method: if $g\ln f\to L$ with $L$ finite, then $f^{g}\to \eu^{L}$; if $g\ln f\to-\infty$ then $f^{g}\to0$; if $g\ln f\to+\infty$ then $f^{g}\to\infty$. Nothing is indeterminate about $g\ln f$ except the product $0\cdot\infty$, and a product is a shape the previous section already knows how to break: divide one factor by the reciprocal of the other, or expand the logarithm.

\begin{keynote}[The transform, once and for all]
Every power-form indeterminacy is the same problem. $1^{\infty}$ gives $\infty\cdot 0$ in the exponent, $\infty^{0}$ gives $0\cdot\infty$, and $0^{0}$ gives $0\cdot(-\infty)$. In all three cases write
\[\lim f^{g}=\exp\Bigl(\lim g\ln f\Bigr),\]
and solve the product. The only hypotheses are that $f>0$ on a punctured neighbourhood of the point and that the inner limit exists in the extended reals. Nothing else about $f$ or $g$ is needed, and in particular neither has to be differentiable.
\end{keynote}

The two forms other than $1^{\infty}$ deserve a moment, because they are easier than they look and are usually dispatched in one line. For $0^{0}$, take $\displaystyle\lim_{x\to0^{+}}x^{x}$: the exponent is $x\ln x\to0$, since $\ln x$ diverges only logarithmically while $x$ vanishes linearly, so the limit is $\eu^{0}=1$. For $\infty^{0}$, take $\displaystyle\lim_{x\to\infty}x^{1/x}$: the exponent is $\tfrac{\ln x}{x}\to0$, so again the limit is $1$. In both cases the decision was made by comparing growth rates inside the product, which is the only skill those two forms require. The reason $1^{\infty}$ gets techniques of its own is that there the product is a genuine contest between two quantities of the same order, and the answer depends on the constant that survives.

The rest of this section is a graded sequence: the general form first, the shortcut that covers most cases second, the memorised $e$-limits third as a corollary of the shortcut rather than as separate facts, and then two structural techniques that use the logarithm for what it does to products rather than to exponents.

\tech[The standard one to the infinity form]{The standard $1^{\infty}$ form}{medium}
\cue The base tends to $1$ and the exponent tends to $\infty$ or to $-\infty$. Typical shapes: $\bigl(\tfrac{\sin x}{x}\bigr)^{1/x^{2}}$, $(\cos x)^{1/x^{2}}$, $\bigl(\tfrac{a^{x}+b^{x}}{2}\bigr)^{1/x}$, and every problem in which a sequence $\bigl(1+\varepsilon_{n}\bigr)^{N_{n}}$ appears with $\varepsilon_{n}\to0$ and $N_{n}\to\infty$.
\method Write $\lim f^{g}=\exp\bigl(\lim g\ln f\bigr)$ and expand the logarithm about $1$: with $u=f-1\to0$,
\[\ln f=\ln(1+u)=u-\frac{u^{2}}{2}+\frac{u^{3}}{3}-\cdots.\]
Keep as many terms of $u$, and as many terms of this expansion, as the exponent's blow-up rate demands: if $g$ grows like $x^{-m}$, you need $\ln f$ correct to order $x^{m}$ inclusive, and no further.

\begin{worked}
\[\lim_{x\to0}\left(\frac{\sin x}{x}\right)^{1/x^{2}}.\]
Here $u=\dfrac{\sin x}{x}-1=-\dfrac{x^{2}}{6}+\dfrac{x^{4}}{120}-\cdots$, so $u^{2}=O(x^{4})$ and the quadratic term of the logarithm cannot reach order $x^{2}$:
\[\ln\frac{\sin x}{x}=u-\frac{u^{2}}{2}+\cdots=-\frac{x^{2}}{6}+O(x^{4}).\]
Multiplying by $g=1/x^{2}$ gives $-\tfrac16+O(x^{2})\to-\tfrac16$, so the limit is $\eu^{-1/6}\approx0.8464817$, and direct evaluation at $x=10^{-3}$ returns $0.8464817$.

Two siblings, computed the same way and worth holding in memory as a set:
\[\lim_{x\to0}\bigl(\cos x\bigr)^{1/x^{2}}=\eu^{-1/2}\approx0.606531,\qquad \lim_{x\to0}\left(\frac{\tan x}{x}\right)^{1/x^{2}}=\eu^{1/3}\approx1.395612,\]
since $\ln\cos x=-\tfrac{x^{2}}{2}+O(x^{4})$ and $\ln\tfrac{\tan x}{x}=\tfrac{x^{2}}{3}+O(x^{4})$. In every case the answer is $\eu^{c}$ where $c$ is simply the coefficient of $x^{2}$ in $\ln f$. There is no pattern connecting the three constants; each has to be computed from its own expansion, which is precisely why the method is a method and not a table.
\end{worked}

The exponent decides how much work the expansion has to do, and this is the only genuine judgement call in the technique. With $g=1/x^{2}$ only the $x^{2}$ coefficient of $\ln f$ survives, and everything else is either infinite (there is nothing of lower order, so the limit is finite) or vanishing. With $g=1/x^{4}$ the $x^{4}$ coefficient becomes the answer, and any surviving $x^{2}$ coefficient becomes an infinity that decides the limit to be $0$ or $\infty$ before the $x^{4}$ term is even reached. The tidy case, and the one problem-setters like, is when the $x^{2}$ term is absent by design.

\begin{keynote}[A fourth-order example worth memorising]
\[\lim_{x\to0}\left(\frac{\cosh x+\cos x}{2}\right)^{1/x^{4}}=\eu^{1/24}\approx1.0425469 .\]
The $x^{2}$ terms of $\cosh x=1+\tfrac{x^{2}}{2}+\tfrac{x^{4}}{24}+\cdots$ and $\cos x=1-\tfrac{x^{2}}{2}+\tfrac{x^{4}}{24}-\cdots$ cancel exactly, leaving the base equal to $1+\tfrac{x^{4}}{24}+O(x^{8})$. Since $u=O(x^{4})$, the term $u^{2}/2$ is $O(x^{8})$ and may be dropped, so $g\ln f\to\tfrac1{24}$. Compare $(\cos x)^{1/x^{4}}$, where the $x^{2}$ term survives, $g\ln f\to-\infty$, and the limit is $0$.
\end{keynote}

\begin{trap}
Stopping at $\ln f\approx f-1$ when the exponent is $1/x^{2}$ or worse and the quadratic term still contributes. Match the expansion order to the exponent: with $g=x^{-m}$, an error of size $x^{m}$ in $\ln f$ is an error of size $1$ in the answer's logarithm, which is a whole factor of $\eu$ in the answer. The symptom is an answer that is off by a clean exponential factor rather than by a little.
\end{trap}

\begin{seealsob}The exponential shortcut without a logarithm; the $e$ limit family; log-transform a product or a nested power.\end{seealsob}

In practice the logarithm is often unnecessary, because to first order $\ln(1+u)$ and $u$ agree. It is worth knowing exactly when the shortcut is legitimate, and the honest statement is both simpler and stronger than the one usually quoted.

\tech[Exponential shortcut without a logarithm]{The $\exp\bigl(\lim g\,(f-1)\bigr)$ shortcut}{medium}
\cue A clean $1^{\infty}$ in which $g\,(f-1)$ visibly tends to a finite limit. This is the common case, and recognising it saves you from writing a logarithm you would only expand again.
\method If $f\to1$, $g\to\infty$ and $g\,(f-1)\to L$ with $L$ finite, then
\[\lim f^{g}=\exp\Bigl(\lim g\,(f-1)\Bigr)=\eu^{L}.\]
The proof is one line. Write $\ln f=(f-1)-\tfrac12(f-1)^{2}\bigl(1+o(1)\bigr)$; then
\[g\,(f-1)^{2}=\frac{\bigl(g(f-1)\bigr)^{2}}{g}\longrightarrow\frac{L^{2}}{\infty}=0,\]
so $g\ln f$ and $g(f-1)$ have the same limit.

\begin{worked}
\[\lim_{x\to0}\bigl(1+\sin2x\bigr)^{\cot3x}.\]
Here $f-1=\sin 2x\to0$ and $g=\cot 3x\to\infty$, so the form is $1^{\infty}$ and the shortcut applies once the product is shown to converge:
\[g\,(f-1)=\frac{\cos3x\,\sin2x}{\sin3x}=\cos 3x\cdot\frac{\sin2x}{2x}\cdot\frac{3x}{\sin3x}\cdot\frac{2}{3}\longrightarrow 1\cdot1\cdot1\cdot\frac23=\frac23 .\]
The limit is $\eu^{2/3}\approx1.9477340$; numerical evaluation at $x=10^{-6}$ gives $1.9477327$.
\end{worked}

\begin{keynote}[The condition is automatic, and that is the point]
Textbooks state the hypothesis as $g\,(f-1)^{2}\to0$, which looks like a condition to be checked. It is not. Whenever $g\to\infty$ and $g(f-1)$ converges to a finite $L$, the quantity $g(f-1)^{2}=\bigl(g(f-1)\bigr)^{2}/g$ tends to $0$ all by itself. The shortcut is therefore valid exactly when it produces a finite answer, and the only real danger is the opposite case, in which $g(f-1)$ diverges and the discarded term is no longer negligible.
\end{keynote}

\begin{trap}
Using the shortcut when $g\,(f-1)\to\pm\infty$ and you need more than the sign. Take $f=1+\tfrac1n$ and $g=n^{2}$. The shortcut reports $\exp(n)$, but the truth is
\[\Bigl(1+\tfrac1n\Bigr)^{n^{2}}=\exp\Bigl(n^{2}\bigl(\tfrac1n-\tfrac1{2n^{2}}+\tfrac{1}{3n^{3}}-\cdots\bigr)\Bigr)=\exp\Bigl(n-\tfrac12+O(\tfrac1n)\Bigr),\]
so $\bigl(1+\tfrac1n\bigr)^{n^{2}}\eu^{-n}\to \eu^{-1/2}\approx0.60653$ and not $1$; at $n=10^{6}$ the ratio is $0.6065309$. The dropped $-\tfrac12$ is a genuine factor. When the exponent diverges, expand the logarithm.
\end{trap}

\begin{seealsob}The standard $1^{\infty}$ form; the $e$ limit family; exponent simplification before base analysis.\end{seealsob}

Everything usually presented as ``the $e$ limits'' is a special case of the shortcut, and it is worth watching the list collapse, because the collapse tells you which of the many memorised forms is the one you actually need. There is only one: the product $g\,(f-1)$.

\tech[The e limit family]{The $e=\lim(1+1/n)^{n}$ family}{easy}
\cue Shapes such as $\bigl(1+\tfrac an\bigr)^{bn}$, $\bigl(\tfrac{n+a}{n+b}\bigr)^{n}$, $(1+ax)^{b/x}$, or any base written as $1+(\text{small})$ with an exponent that is the reciprocal of that small quantity, up to a constant factor.
\method Apply the shortcut. The master forms are
\[\lim_{n\to\infty}\Bigl(1+\frac an\Bigr)^{bn}=\eu^{ab},\qquad \lim_{x\to0}(1+ax)^{b/x}=\eu^{ab},\]
both because $g\,(f-1)=bn\cdot\tfrac an=ab$ identically, with no expansion at all. For $\bigl(\tfrac{n+a}{n+b}\bigr)^{n}$, rewrite the base as $1+\tfrac{a-b}{n+b}$ and read off $\eu^{a-b}$; for example $\bigl(\tfrac{n+3}{n+1}\bigr)^{n}\to\eu^{2}\approx7.389056$. When the deviation is of order $n^{-2}$ the exponent must be of order $n^{2}$ to match: $\bigl(\tfrac{n^{2}+1}{n^{2}-2}\bigr)^{n^{2}}=\bigl(1+\tfrac{3}{n^{2}-2}\bigr)^{n^{2}}\to\eu^{3}\approx20.0855$, which at $n=10^{3}$ evaluates to $20.0864$.

\begin{worked}
\[\lim_{x\to0}\left(\frac{a^{x}+b^{x}}{2}\right)^{1/x},\qquad a,b>0.\]
The base tends to $\tfrac{1+1}{2}=1$, so the form is $1^{\infty}$, and
\[g\,(f-1)=\frac1x\cdot\frac{(a^{x}-1)+(b^{x}-1)}{2}=\frac12\left(\frac{a^{x}-1}{x}+\frac{b^{x}-1}{x}\right)\longrightarrow\frac{\ln a+\ln b}{2}=\ln\sqrt{ab},\]
using the standard limit $\tfrac{a^{x}-1}{x}\to\ln a$. The limit is therefore $\sqrt{ab}$: the arithmetic mean of the exponentials produces the geometric mean of the bases. At $a=2$, $b=5$ this is $\sqrt{10}\approx3.162278$, which is what numerical evaluation at $x=10^{-6}$ returns.
\end{worked}

\begin{trap}
Mismatching the two rates. In $\bigl(1+\tfrac{a}{n^{2}}\bigr)^{n}$ the product $g(f-1)=a/n\to0$, so the limit is $1$, not $\eu^{a}$. The exponent must be the reciprocal of the base's deviation from $1$ to within a constant factor; if the exponent is smaller in order the answer is $1$, and if it is larger the answer is $0$ or $\infty$. Always compute the product rather than pattern-matching the shape.
\end{trap}

\begin{seealsob}The standard $1^{\infty}$ form; the exponential shortcut without a logarithm; the $n$th root of a sum tends to the maximum.\end{seealsob}

Before invoking any of this machinery, look at the exponent on its own. A surprising number of frightening powers are not indeterminate at all once the exponent has been simplified, and running the logarithm on a determinate form wastes half a page and invites the sign errors that come with it.

\tech{Exponent simplification before base analysis}{medium}
\cue A power whose \emph{exponent} is itself an indeterminate quotient, typically at $x\to1$ with radicals present: exponents like $\tfrac{1-x}{1-\sqrt x}$, $\tfrac{\ln x}{x-1}$, or $\tfrac{x^{2}-1}{x-1}$.
\method Resolve the exponent first, using the surgery of the previous section. If it tends to a finite constant $c$ and the base tends to a positive $B\neq1$, the whole limit is simply $B^{c}$, by continuity of $(b,t)\mapsto b^{t}$ on $(0,\infty)\times\R$. Only if the base also tends to $1$ do you need the logarithm, and by then the exponent is a known constant and the remaining work is halved.

\begin{worked}
\[\lim_{x\to1}\left(\frac{x}{1+x^{2}}\right)^{\frac{1-x}{1-\sqrt{x}}}.\]
The exponent factors: $1-x=(1-\sqrt x)(1+\sqrt x)$, so $\dfrac{1-x}{1-\sqrt x}=1+\sqrt x\to2$. The base tends to $\tfrac{1}{2}$, which is neither $0$ nor $1$, so the form is determinate and the answer is $\bigl(\tfrac12\bigr)^{2}=\tfrac14$, confirmed numerically at $x=1-10^{-7}$.

Change the base to $\dfrac{2x}{1+x^{2}}$, which tends to $1$, and the same exponent analysis leaves a genuine $1^{\infty}$; here the base is $1-\tfrac{(1-x)^{2}}{1+x^{2}}$, so $g\,(f-1)\to 2\cdot0=0$ and the limit is $1$. The point is that in both cases the exponent was settled before any decision about machinery was made.
\end{worked}

\begin{trap}
Logarithmising a form that is not indeterminate. If the base tends to $\tfrac12$ and the exponent to $2$, then $\ln f\to-\ln 2$ and $g\ln f\to-2\ln2$: correct, but three unnecessary lines, and along the way it is easy to write $\lim g\ln f=\lim g\cdot\lim\ln f$ in a situation where that product rule genuinely does not apply, because the exponent's own limit had not yet been established.
\end{trap}

\begin{seealsob}Change of variable to move the limit point; the standard $1^{\infty}$ form; factor out the removable singularity.\end{seealsob}

The logarithm does more than tame exponents. Its structural role is that it converts multiplication into addition, and sums are the objects for which we have the largest arsenal: telescoping, comparison, integral bounds, and (in the next chapter) recognition as integrals. Any time a limit contains a product of many factors, the first move is to take its logarithm and see what kind of sum appears.

\tech{Log-transform a product or a nested power}{hard}
\cue An infinite product, an $n$th root of a product of $n$ factors, a tower of powers, or a sum of logarithms that would be a product if you read it backwards.
\method Use $\ln\prod_{k}c_{k}=\sum_{k}\ln c_{k}$. Roots of products become averages of logarithms, infinite products become series, and the whole apparatus for sums becomes available. Exponentiate only at the very end. Before taking logarithms, confirm that the factors are eventually positive, and if finitely many are not, split them off as a separate constant factor.

\begin{worked}
A finite warm-up in which the product telescopes and no analysis is needed:
\[\prod_{k=2}^{n}\Bigl(1-\frac{1}{k^{2}}\Bigr)=\prod_{k=2}^{n}\frac{(k-1)(k+1)}{k^{2}}=\frac{1}{2}\cdot\frac{n+1}{n}\longrightarrow\frac12 ,\]
where the numerators and denominators cancel in two interleaved chains.

A second finite product, telescoping through a trigonometric identity rather than an arithmetic one. Repeated use of $\sin2\theta=2\sin\theta\cos\theta$ gives
\[\prod_{k=1}^{n}\cos\frac{x}{2^{k}}=\frac{\sin x}{2^{n}\sin\bigl(x/2^{n}\bigr)}\longrightarrow\frac{\sin x}{x}\quad(n\to\infty),\]
because $2^{n}\sin(x/2^{n})\to x$. At $x=1$ the partial products agree with $\sin1=0.8414710$ to fifteen digits. This is Vi\`ete's product, and it is the cleanest example in the book of a product whose logarithm you should \emph{not} take: the telescoping is visible in the product itself.

Now the serious case. Euler's product $\dfrac{\sin x}{x}=\displaystyle\prod_{n\geq1}\Bigl(1-\dfrac{x^{2}}{n^{2}\pi^{2}}\Bigr)$ converges for every $x$, and for $0<\abs{x}<\pi$ every factor is positive, so taking logarithms is legitimate:
\[\sum_{n=1}^{\infty}\ln\Bigl(1-\frac{x^{2}}{n^{2}\pi^{2}}\Bigr)=\ln\frac{\sin x}{x}=-\frac{x^{2}}{6}+O(x^{4}).\]
Dividing by $x^{2}$,
\[\lim_{x\to0}\frac{1}{x^{2}}\sum_{n=1}^{\infty}\ln\Bigl(1-\frac{x^{2}}{n^{2}\pi^{2}}\Bigr)=-\frac16 ,\]
and evaluating the series at $x=10^{-3}$ gives $-0.1666667$. Reading the same identity termwise, $\ln(1-t)=-t+O(t^{2})$ turns the sum into $-\tfrac{x^{2}}{\pi^{2}}\sum n^{-2}$, so the answer is $-\tfrac16$ precisely because $\sum n^{-2}=\pi^{2}/6$. The two routes agree, which is Euler's original argument run backwards.
\end{worked}

\begin{trap}
Exchanging $\ln$ with an infinite sum, or a limit with an infinite sum, without a convergence argument. The termwise reading above is legitimate because $\sum x^{2}/(n^{2}\pi^{2})$ converges uniformly on compact subsets of $(-\pi,\pi)$. The second failure is silent: $\ln$ of a negative factor is undefined, so a product whose early factors are negative must have them split off before the logarithm is applied, and a product with a zero factor has no logarithm at all.
\end{trap}

\begin{seealsob}Riemann sum via the logarithm of a product; the standard $1^{\infty}$ form; nested limits taken strictly in order.\end{seealsob}

The last technique of this section is a discipline rather than a manipulation, and it becomes unavoidable the moment a problem contains two limit operators. It is placed here because the commonest source of a hidden double limit is precisely an inner $1^{\infty}$ whose value is a function of an outer variable, so the two halves of the problem use the two halves of this section.

\tech{Nested limits, taken strictly in order}{hard}
\cue Two limit symbols, an inner limit that defines a function of the outer variable, or a sequence of functions whose limit is then evaluated at a moving point.
\method Complete the inner limit \emph{fully}, obtaining a closed form in the outer variable, and only then take the outer limit. The two orders generally disagree, and neither may be interchanged without a uniformity hypothesis such as uniform convergence or monotonicity.

\begin{worked}
Let $g(x)=\displaystyle\lim_{r\to0}\Bigl((x+1)^{r+1}-x^{r+1}\Bigr)^{1/r}$ for $x>0$, and evaluate $\displaystyle\lim_{x\to\infty}\frac{g(x)}{x}$.

\emph{Inner limit.} Put $h(r)=(x+1)^{r+1}-x^{r+1}$, so $h(0)=(x+1)-x=1$: the inner form is $1^{\infty}$ in $r$. Because $\ln h(0)=0$, the quotient $\tfrac{\ln h(r)}{r}$ is itself a difference quotient, and
\[\lim_{r\to0}\frac{\ln h(r)}{r}=\bigl(\ln h\bigr)'(0)=\frac{h'(0)}{h(0)}=h'(0)=(x+1)\ln(x+1)-x\ln x,\]
since $h'(r)=(x+1)^{r+1}\ln(x+1)-x^{r+1}\ln x$. Exponentiating,
\[g(x)=\exp\bigl((x+1)\ln(x+1)-x\ln x\bigr)=\frac{(x+1)^{x+1}}{x^{x}} .\]
\emph{Outer limit.} With the closed form in hand the second limit is routine:
\[\frac{g(x)}{x}=\frac{(x+1)^{x+1}}{x^{x+1}}=\Bigl(1+\frac1x\Bigr)^{x+1}\longrightarrow \eu ,\]
by the $e$ limit family. Numerically $g(10^{5})/10^{5}=2.71859$ against $\eu=2.71828$.
\end{worked}

\begin{trap}
Collapsing both limits at once, or swapping them. The standard warning is
\[\lim_{m\to\infty}\lim_{n\to\infty}\frac{m}{m+n}=0,\qquad \lim_{n\to\infty}\lim_{m\to\infty}\frac{m}{m+n}=1 .\]
In the worked example, letting $r\to0$ and $x\to\infty$ together loses the fact that the inner answer depends on $x$ through $(x+1)\ln(x+1)-x\ln x$, and that dependence is exactly what produces $\eu$ rather than $1$. A sharper pair, where both iterated limits exist and disagree by more than a constant: with $a_{m,n}=\bigl(1+\tfrac1m\bigr)^{n}$, the inner limit over $n$ is $\infty$ for every fixed $m$, while the inner limit over $m$ is $1$ for every fixed $n$. Neither order is wrong; they are answers to different questions.
\end{trap}

\begin{seealsob}Log-transform a product or a nested power; the standard $1^{\infty}$ form; iterated versus joint limits.\end{seealsob}

One closing observation ties the section together. Every technique in it was an application of a single sentence: a power is an exponential of a product, and $\exp$ is continuous. The generalisation of the geometric-mean example is immediate and worth doing yourself once,
\[\lim_{x\to0}\left(\frac{a_{1}^{x}+a_{2}^{x}+\cdots+a_{m}^{x}}{m}\right)^{1/x}=\bigl(a_{1}a_{2}\cdots a_{m}\bigr)^{1/m},\]
for positive $a_{i}$, because $g(f-1)=\tfrac1m\sum_{i}\tfrac{a_{i}^{x}-1}{x}\to\tfrac1m\sum_{i}\ln a_{i}$. At $m=3$ with $(2,3,7)$ the value is $\sqrt[3]{42}\approx3.476027$, matched by direct evaluation at $x=10^{-6}$. Nothing new was needed for the general case, which is the sign that the machine, and not the individual answers, is the thing to learn.

\section{The squeeze}

The first two sections both assume that the expression can be transformed into one whose parts converge. Some expressions cannot. If $f(x)=x^{2}\sin(1/x)$, then $\sin(1/x)$ has no limit at $0$ and no rewriting will give it one; the limit of $f$ nevertheless exists, and the only way to see it is to trap $f$ between two functions that do converge. This is the squeeze, and it is the tool of last resort in exactly the sense that it is the tool that never needs the expression to cooperate.

The squeeze also has a second life as a technique for sums whose length grows with the limit variable, where it competes directly with the Riemann-sum recognition of the next chapter. The rule of thumb is worth stating in advance, because it decides which chapter a problem belongs to: if the crude bounds, meaning the smallest term times the count and the largest term times the count, already agree in the limit, use the squeeze and stop. If they disagree, the sum has genuine internal structure across its range and you need the integral.

\tech{The squeeze theorem template}{easy}
\cue An expression that resists exact evaluation but is visibly trapped between two computable bounds. Oscillating factors, floor and fractional-part functions, and anything defined by an inequality rather than a formula are the standard signatures.
\method Produce $\ell(x)\leq f(x)\leq u(x)$ on a punctured neighbourhood of $a$, show $\lim_{x\to a}\ell=\lim_{x\to a}u=L$, and conclude $\lim_{x\to a}f=L$. The bounds may be crude; crudeness costs nothing provided both ends still converge to the same value. The theorem holds verbatim for sequences, for one-sided limits, and for $a=\pm\infty$.

\begin{worked}
\[\lim_{x\to0^{+}}x\floorb{\frac1x}.\]
The defining inequality of the floor is $\tfrac1x-1<\floorb{\tfrac1x}\leq\tfrac1x$ for every $x>0$. Multiplying through by the positive number $x$ preserves both directions:
\[1-x<x\floorb{\frac1x}\leq1 .\]
Both ends tend to $1$, so the limit is $1$, even though $\floorb{1/x}$ jumps infinitely often in every neighbourhood of $0$ and the function is discontinuous at infinitely many points of any interval $(0,\delta)$.

The baby case, and the one to reach for when an oscillation is bounded: $-x^{2}\leq x^{2}\sin\tfrac1x\leq x^{2}$ gives $\displaystyle\lim_{x\to0}x^{2}\sin\tfrac1x=0$, with no claim whatsoever about $\sin(1/x)$ itself.

A sequence version, where the bound is built by comparing consecutive terms. For fixed $a>0$ choose an integer $m>2a$. For $n>m$ every further factor multiplies by $a/k<\tfrac12$, so
\[0<\frac{a^{n}}{n!}\leq\frac{a^{m}}{m!}\Bigl(\frac12\Bigr)^{n-m}\longrightarrow0,\]
and therefore $a^{n}/n!\to0$ for every $a$. At $a=20$, $n=300$ the value is about $6.7\times10^{-225}$, though the sequence increases until $n=20$ before it turns. The squeeze is indifferent to that early growth, which is exactly why it is the right tool.
\end{worked}

\begin{trap}
Bounds that converge to different values prove nothing at all, and bounds valid only on one side of $a$ prove only a one-sided limit. The multiplication step above is the other classic failure: multiplying an inequality by $x$ reverses it when $x<0$, so the corresponding limit as $x\to0^{-}$ needs its own argument with the inequalities written the other way round.
\end{trap}

\begin{seealsob}Bounded times infinitesimal; squeeze on a sum whose length is the limit variable; integral bounds on a sum.\end{seealsob}

\begin{keynote}[The squeeze proves the facts everything else assumes]
$\displaystyle\lim_{n\to\infty}n^{1/n}=1$. Write $n^{1/n}=1+h_{n}$ with $h_{n}>0$ for $n\geq2$. The binomial theorem gives $n=(1+h_{n})^{n}\geq\binom{n}{2}h_{n}^{2}=\tfrac{n(n-1)}{2}h_{n}^{2}$, so $0<h_{n}\leq\sqrt{\tfrac{2}{n-1}}\to0$. There is no way to obtain this by algebra on $n^{1/n}$ itself, and every use of the root test later in the book depends on it. The same argument gives $a^{1/n}\to1$ for every $a>0$.
\end{keynote}

The single most frequent application of the squeeze is so common that it deserves its own name and its own reflex, because the alternative reasoning that students reach for is not merely inefficient but invalid.

\begin{keynote}[The squeeze run backwards proves nonexistence]
The squeeze establishes a limit by trapping; the same idea run backwards refutes one. If two sequences $x_{n}\to a$ and $y_{n}\to a$ (with $x_{n},y_{n}\neq a$) give $f(x_{n})\to L_{1}$ and $f(y_{n})\to L_{2}$ with $L_{1}\neq L_{2}$, then $\lim_{x\to a}f(x)$ does not exist. For $f(x)=\sin(1/x)$ at $a=0$, take $x_{n}=\tfrac{1}{2\pi n}$ and $y_{n}=\tfrac{1}{\pi/2+2\pi n}$: the values are $0$ and $1$ for every $n$. This is why $x^{2}\sin(1/x)\to0$ has to be proved by bounding and cannot be proved by a product rule: the second factor genuinely has no limit, and the sequence argument shows it in two lines.
\end{keynote}

\tech{Bounded times infinitesimal}{easy}
\cue A product $b(x)\,\varepsilon(x)$ where $b$ is bounded, typically an oscillating $\sin$, $\cos$, fractional part, or $\sgn$, and $\varepsilon\to0$. Also a quotient in which such a bounded term is added to something that dominates it.
\method If $\abs{b}\leq M$ on a punctured neighbourhood and $\varepsilon\to0$, then $b\varepsilon\to0$. The proof is $0\leq\abs{b\varepsilon}\leq M\abs{\varepsilon}\to0$ followed by the squeeze applied to $\abs{b\varepsilon}$. Crucially, no assumption is made that $\lim b$ exists, and none is available in the interesting cases.

\begin{worked}
\[\lim_{x\to\infty}\frac{x+\sin x}{x+\cos x}=\lim_{x\to\infty}\frac{1+\dfrac{\sin x}{x}}{1+\dfrac{\cos x}{x}}=\frac{1+0}{1+0}=1 .\]
Dividing through by $x$ is what converts the two non-convergent terms into bounded-times-infinitesimal products; each tends to $0$ by the bound $\abs{\sin x}\leq1$, and only after that step is the quotient rule for limits legal, because only then do numerator and denominator separately have limits.

The same bound settles fractional parts, where the oscillation is not even continuous. Writing $\{t\}=t-\floorb{t}$ for the fractional part, $0\leq\{1/x\}<1$ for all $x\neq0$, so
\[\lim_{x\to0}x\Bigl\{\frac1x\Bigr\}=0\quad\text{and}\quad \lim_{x\to0}\Bigl(x\Bigl\{\frac1x\Bigr\}+x\floorb{\tfrac1x}\Bigr)=\lim_{x\to0}x\cdot\frac1x=1,\]
the second because the two pieces reassemble into something exact. The first limit is bounded times infinitesimal; the second is a reminder that the decomposition, not the individual pieces, is what carries the information.
\end{worked}

\begin{trap}
Writing $\displaystyle\lim_{x\to\infty}\frac{\sin x}{x}=\lim_{x\to\infty}\sin x\cdot\lim_{x\to\infty}\frac1x$. The product rule for limits requires both factors to have limits; here the first has none, and the correct justification is the boundedness estimate. The conclusion happens to be right, which is what makes the habit dangerous: the same reasoning applied to $\dfrac{x\sin x}{x}$ would produce $0$ instead of the correct answer, which is that the limit does not exist.
\end{trap}

\begin{seealsob}The squeeze theorem template; squeeze on a sum whose length is the limit variable; the $n$th root of a sum tends to the maximum.\end{seealsob}

Both of the last two blocks share a feature worth naming: the bound was found by discarding information deliberately. Replacing $\sin x$ by $1$, or $\floorb{1/x}$ by $1/x$, throws away everything that made the expression hard, and the limit survives the loss. That is the whole art. A bound that keeps too much detail is as useless as no bound, because it will be as hard to evaluate as the original; a bound that throws away too much will not pinch. Aim first for the crudest bound imaginable, check whether both ends agree, and only sharpen if they do not.

Sums whose number of terms grows are the second natural home of the squeeze. Every individual term tends to $0$, so the termwise reading gives $0$; the sum does not. Bounding all the terms by the two extreme ones resolves the tension in a single step, whenever it works at all, and the diagnostic for whether it works is worth internalising because it also tells you when to turn the page to the Riemann sums.

\tech{Squeeze on a sum whose length is the limit variable}{medium}
\cue $S_{n}=\sum_{k=1}^{n}a_{n,k}$ where the number of terms grows with $n$ and each individual term tends to $0$. The tell is that the terms vary only slightly across $k$, so that the first and last are nearly equal.
\method Bound every summand by the extremes:
\[n\cdot\min_{1\leq k\leq n}a_{n,k}\;\leq\;S_{n}\;\leq\;n\cdot\max_{1\leq k\leq n}a_{n,k}.\]
If both bounds share a limit you are finished, and no integral is needed. If they do not, the sum's internal variation is of the same order as the sum itself, and the Riemann-sum machinery of the next chapter is required instead.

\begin{worked}
\[\lim_{n\to\infty}\sum_{k=1}^{n}\frac{1}{\sqrt{n^{2}+k}}.\]
The terms decrease in $k$, so the largest is at $k=1$ and the smallest at $k=n$:
\[\frac{n}{\sqrt{n^{2}+n}}\leq S_{n}\leq\frac{n}{\sqrt{n^{2}+1}} .\]
Both ends tend to $1$, so the limit is $1$; at $n=10^{6}$ the sum is $0.9999997$. The reason the squeeze pinches is that $k\leq n$ is negligible beside $n^{2}$, so the ratio of the extreme terms tends to $1$. The same test passes for $\displaystyle\sum_{k=1}^{n}\frac{k}{n^{2}+k}$, whose extreme terms are $\tfrac{1}{n^{2}+1}$ and $\tfrac{n}{n^{2}+n}$: here the crude bounds do \emph{not} pinch, but the sum is $\tfrac{1}{n^{2}}\sum k$ to within $O(1/n)$, so the limit is $\tfrac12$, and at $n=10^{6}$ the value is $0.5000002$. Read that as a warning: failing the ratio test for the crude bounds means only that these particular bounds fail, not that the sum is hard.

Contrast $\displaystyle\sum_{k=1}^{n}\frac{1}{\sqrt{n^{2}+k^{2}}}$, where the same bounds give only $\tfrac{1}{\sqrt2}\leq S_{n}\leq1$ and refuse to close. Here $k$ enters at the same scale as $n$, the extreme terms differ by a factor of $\sqrt2$ in the limit, and the answer is the integral
\[\int_{0}^{1}\frac{\dd u}{\sqrt{1+u^{2}}}=\ln\bigl(1+\sqrt2\bigr)\approx0.881374 .\]
\end{worked}

\begin{trap}
Taking the limit termwise and concluding $0+0+\cdots=0$. The sum rule for limits extends to a \emph{fixed} number of summands only, and here the number of summands is the limit variable itself. The diagnostic for which tool applies is the ratio of the largest term to the smallest: if it tends to $1$, the squeeze pinches; if it tends to a constant greater than $1$, it cannot, and no refinement of the same two bounds will help.
\end{trap}

\begin{seealsob}Riemann sum recognition; the squeeze theorem template; integral bounds on a sum.\end{seealsob}

When the summand is monotone in $k$ but its variation across the range is too large for the crude bounds, there is still a squeeze available, with the two bounds supplied by integrals rather than by extreme terms. This is a sharper instrument, it is the source of most discrete asymptotics, and it is the bridge to the next chapter: the same picture of rectangles against a curve reappears there with the roles of sum and integral exchanged.

\tech{Integral bounds on a sum}{medium}
\cue A sum of the values of a monotone function, where you need the asymptotic size rather than the exact value: $\sum k^{p}$, $\sum 1/k$, $\sum\ln k$, or a tail $\sum_{k>n}f(k)$ whose size controls a remainder.
\method For $f$ decreasing on $[1,\infty)$,
\[\int_{1}^{n+1}f(x)\dd x\;\leq\;\sum_{k=1}^{n}f(k)\;\leq\;f(1)+\int_{1}^{n}f(x)\dd x,\]
and for $f$ increasing the inequalities reverse in the obvious way, becoming $\int_{0}^{n}f\leq\sum_{k=1}^{n}f(k)\leq\int_{1}^{n+1}f$. This is simultaneously the integral test for convergence and the source of $H_{n}=\ln n+\gamma+O(1/n)$.

\begin{worked}
\[\lim_{n\to\infty}\frac{1}{n^{p+1}}\sum_{k=1}^{n}k^{p},\qquad p>0.\]
Since $x^{p}$ increases, the rectangle of height $k^{p}$ and width $1$ sits above $\int_{k-1}^{k}x^{p}\dd x$ and below $\int_{k}^{k+1}x^{p}\dd x$, so summing,
\[\frac{n^{p+1}}{p+1}=\int_{0}^{n}x^{p}\dd x\;\leq\;\sum_{k=1}^{n}k^{p}\;\leq\;\int_{1}^{n+1}x^{p}\dd x=\frac{(n+1)^{p+1}-1}{p+1}.\]
Dividing by $n^{p+1}$, both ends tend to $\dfrac{1}{p+1}$. At $p=3$ the limit is $\tfrac14$, and the sum at $n=2\times10^{5}$ gives $0.2500025$.

The same inequality with $f(x)=1/x$, now decreasing, gives $\ln(n+1)\leq H_{n}\leq1+\ln n$, so $H_{n}-\ln n$ is trapped in $[0,1]$; the sequence is decreasing and its limit is $\gamma\approx0.5772157$. At $n=10^{6}$ the difference is $0.5772162$. With $f(x)=\ln x$ the same picture gives $\ln(n!)=n\ln n-n+O(\ln n)$, which is Stirling's formula without its constant; at $n=10^{6}$ the discrepancy $\ln(n!)-(n\ln n-n)$ is $7.8267$, against $\tfrac12\ln(2\pi n)=7.8267$.

The same inequality run on a tail rather than a head is how remainders are estimated. For $f$ decreasing and positive,
\[\int_{n+1}^{\infty}f\;\leq\;\sum_{k>n}f(k)\;\leq\;\int_{n}^{\infty}f,\]
so for $f(x)=x^{-2}$ the tail $\sum_{k>n}k^{-2}$ is trapped between $\tfrac{1}{n+1}$ and $\tfrac1n$, giving $n\sum_{k>n}k^{-2}\to1$. The two bounds differ by a term of order $n^{-2}$, which is what makes the estimate useful rather than merely true.
\end{worked}

\begin{trap}
Using the wrong endpoint convention and losing a whole term. Draw the rectangles: for a decreasing $f$, the rectangle of width $1$ and height $f(k)$ sits \emph{above} $\int_{k}^{k+1}f$ and \emph{below} $\int_{k-1}^{k}f$, and the direction of both inequalities flips when $f$ increases. A misplaced endpoint changes $H_{n}$ by $1$, which is the difference between a convergent and a divergent conclusion in some comparison arguments.
\end{trap}

\begin{seealsob}Squeeze on a sum whose length is the limit variable; Riemann sum recognition; the integral test.\end{seealsob}

The final technique of the chapter is the purest illustration of the squeeze's power: an expression involving several competing exponentials, with no algebraic simplification available and no logarithm that helps, whose limit falls out of two lines of bounding.

\tech[The nth root of a sum tends to the maximum]{The $n$th root of a sum tends to the maximum}{easy}
\cue $\bigl(a_{1}^{n}+\cdots+a_{m}^{n}\bigr)^{1/n}$ with $m$ \emph{fixed} and all $a_{i}\geq0$; equivalently the $L^{p}$ norm of a fixed finite list as $p\to\infty$.
\method Let $M=\max_{i}a_{i}>0$. Then $M^{n}\leq\sum_{i}a_{i}^{n}\leq m\,M^{n}$, since every term is at most $M^{n}$ and at least one equals it. Taking $n$th roots,
\[M\;\leq\;\Bigl(\sum_{i=1}^{m}a_{i}^{n}\Bigr)^{1/n}\;\leq\;M\,m^{1/n}\longrightarrow M,\]
since $m^{1/n}=\eu^{(\ln m)/n}\to1$ for fixed $m$. The squeeze does all the work; the largest base swallows the rest.

\begin{worked}
$\displaystyle\lim_{n\to\infty}\bigl(2^{n}+3^{n}\bigr)^{1/n}=3$, because $3\leq(2^{n}+3^{n})^{1/n}\leq3\cdot2^{1/n}\to3$; at $n=200$ the value already agrees with $3$ to fifteen digits. Equally $\bigl(1^{n}+2^{n}+\cdots+7^{n}\bigr)^{1/n}\to7$, with the crude upper bound $7\cdot7^{1/n}$. Constants in front change nothing: $\Bigl(\tfrac{a^{n}+b^{n}}{2}\Bigr)^{1/n}\to\max(a,b)$, since the factor $2^{-1/n}\to1$ is absorbed by the same argument. Compare this with the geometric mean produced by $\bigl(\tfrac{a^{x}+b^{x}}{2}\bigr)^{1/x}$ as $x\to0$: identical expression, opposite limit point, and the answers are the two classical means.

The continuous analogue is the same argument: for a step function taking finitely many values, $\bigl(\int_{0}^{1}\abs{f}^{p}\bigr)^{1/p}\to\norm{f}_{\infty}$ as $p\to\infty$, with $m$ playing the role of the number of levels and the measure of each level playing the role of the constant that $m^{1/n}$ absorbs.
\end{worked}

\begin{trap}
Letting $m$ grow with $n$. The factor $m^{1/n}$ then need not tend to $1$ and the conclusion fails outright: take $m=2^{n}$ copies of $a_{i}=1$, so the sum is $2^{n}$ and its $n$th root is $2$, while the maximum is $1$. The hypothesis that $m$ is fixed is doing real work, and any problem in which the number of competing terms grows must be handled by the sum techniques instead.
\end{trap}

\begin{seealsob}The squeeze theorem template; the $e$ limit family; the root test for series.\end{seealsob}

Stand back from the three sections and the division of labour is clean. Surgery works when the expression is lying about itself, and it succeeds by producing an equal expression that tells the truth. The exponential machine works when the expression is honest but the operation is wrong, and it succeeds by replacing that operation with one that limits respect. The squeeze works when the expression contains something that will never converge, and it succeeds by refusing to look inside. Nothing in this chapter needed a derivative, a Taylor series, or L'H\^opital's rule, and a solver who reaches for those first will spend three times as long on most of these problems and will still get the structural questions wrong.

\section{Recognition matrix}
\begin{recog}{@{}p{0.31\textwidth}p{0.35\textwidth}p{0.28\textwidth}@{}}
\rowcolor{mist}\mh{If you see} & \mh{Reach for} & \mh{If that fails} \\
\midrule
\endhead
$\sqrt{A}-\sqrt{B}$ in a $\tfrac00$ or $\infty-\infty$ & Conjugate: $\dfrac{A-B}{\sqrt A+\sqrt B}$ & Bridging term at the common limit \\
$\sqrt[3]{A}-\sqrt[3]{B}$ & Cube conjugate $a^{2}+ab+b^{2}$ & Binomial expansion of $(1+u)^{1/3}$ \\
$\sqrt[3]{f}-\sqrt{g}$, both $\to1$ & Bridge at $1$, then two conjugates & Expand both to order $x^{2}$ \\
Polynomial $\tfrac00$ at $x=a$ & Divide out $(x-a)$, possibly twice & Synthetic division, then substitute \\
Two fractions each blowing up & Common denominator first & Change of variable $h=x-a$ \\
Roots of orders $m$ and $n$ at $x=1$ & Substitute $x=t^{\lcm(m,n)}$ & Difference-quotient reading \\
Limit point $\pi/2$, $1$, or $\infty$ & Move it to $0$: $u=x-a$ or $t=1/x$ & Track the direction of approach \\
Base $\to1$, exponent $\to\infty$ & $\exp\bigl(\lim g(f-1)\bigr)$ & $\exp\bigl(\lim g\ln f\bigr)$, $\ln(1+u)$ expanded \\
Base $\to1$, exponent like $x^{-4}$ & Expand $\ln f$ to order $x^{4}$ & Check whether the $x^{2}$ term cancels \\
Exponent is itself a $\tfrac00$ quotient & Simplify the exponent first & Then test whether the base is $1$ \\
$0^{0}$ or $\infty^{0}$ & Same transform: $\exp(g\ln f)$ & Compare growth rates in the product \\
A product, or an $n$th root of one & Take logarithms: $\ln\prod=\sum\ln$ & Read the sum as an integral \\
Two limit operators & Inner limit to closed form, then outer & Check that the orders disagree \\
A bounded but non-convergent factor & Bounded times infinitesimal & Squeeze with explicit $\ell$ and $u$ \\
$\sum_{k=1}^{n}a_{n,k}$, terms nearly equal & $n\min\leq S_{n}\leq n\max$ & Riemann sum, or integral bounds \\
$\sum f(k)$ with $f$ monotone & Integral bounds on the sum & Euler and Maclaurin correction terms \\
$\bigl(a_{1}^{n}+\cdots+a_{m}^{n}\bigr)^{1/n}$, $m$ fixed & Factor out the maximum & Check that $m$ really is fixed \\
\end{recog}
